% !TEX root = ../paper.tex

\section{Introduction}
\label{sec:intro}

Housing costs in the United States have risen sharply, and a large literature points to land use regulation as a central cause.
Restrictive zoning, minimum lot sizes, and density caps limit housing supply, raise prices, and distort the allocation of labor across cities \citep{glaeser_why_2005, saiz_geographic_2010, turner_land_2014, hsieh_housing_2019, duranton_urban_2023}.
Formal rules, however, are only one part of the regulatory environment.
A small but growing body of work argues that the political environment and fragmentation of regulatory authority matter as much as the formal regulations themselves \citep{mayer_land_2000, hilber_origins_2013, khan_decentralized_2021, mast_warding_2024, soltas_how_2026}.
When developers face long and unpredictable timelines, projects shrink, marginal investments become unviable, and housing supply falls \citep{cunningham_house_2006, bloom_impact_2009, oh_land_2024}.
Despite this preliminary evidence, there is little empirical work on how individual politicians use these informal regulatory powers to shape housing supply.

This paper studies this informal regulatory power in the context of Chicago's ``aldermanic privilege,'' the unwritten but influential power of each of Chicago's 50 city council members (aldermen) to unilaterally block or approve development in their ward.
Chicago has no city charter, so there is no statutory basis for this power.
It operates through deference: other aldermen vote with the local member on land use matters, city agencies follow the local alderman's lead, and developers understand that major projects require the alderman's blessing \citep{thale_aldermanic_2005, einstein_neighborhood_2019, khan_decentralized_2021}.
This creates sharp changes in regulatory authority at ward boundaries, even when the neighborhood and zoning are quite similar.
Two parcels on opposite sides of the same street, in the same neighborhood, with the same zoning, can face very different regulatory regimes depending on which alderman they interact with.
I use this sharp change in regulatory authority to study how informal regulations affect housing supply and prices.

To measure the regulatory environment each alderman creates, I construct ``regulatory stringency scores'' for each alderman from permit processing times for projects where aldermen are most likely to matter: new construction, renovation, demolition, porch construction, and reinstated permits.
I residualize processing times against ward characteristics and time effects to isolate the alderman-specific component of the variation, then apply empirical Bayes shrinkage \citep{kline_systemic_2022, kline_discrimination_2024} to reduce noise for aldermen with relatively few permits.
The motivation for using processing time variation is simple: delay makes projects less valueable, and
a developer who expects years of meetings, redesigns, and uncertainty may build a smaller project or be deterred from building at all \citep{kydland_time_1982, mayer_land_2000, paciorek_supply_2013, gyourko_local_2021, soltas_how_2026}.

I then estimate effects on housing supply by comparing new construction density near ward boundaries.
On one side, developers deal with one alderman; across the street, they deal with another, which can result in a completely different regulatory environment.
The neighborhood and statuatory zoning are similar, but the informal regulatory environment changes sharply.
I use a spatial regression discontinuity design to compare new construction projects on either side of ward boundaries, and find that moving to a more stringent alderman substantially reduces the density of new construction, particularly for multifamily projects.
For multifamily new construction, floor-area ratio (FAR) is about $25\%$ lower and dwelling units per acre (DUPAC) is about $33\%$ lower on the more stringent side of the boundary.
The all-construction estimates point in the same direction but are smaller, at about $11\%$ lower FAR and $14\%$ lower DUPAC.

I then complement this analysis by using the 2015 ward redistricting to test whether permitting changes when a block receives a new alderman, keeping location fixed.
Using this within-census-block variation, I find that blocks moved to more stringent aldermen experience a 7.4\% decline for these `high-discretion' permits over the next 5 years, relative to nearby blocks that remain with the same alderman and thus see no change in the regulatory environment.
This result shows that more stringent aldermen reduce not only the scale of completed projects, but also the flow of projects requiring discretionary review, suggesting that stringent aldermen deter development on the extensive margin as well.

I finally examine whether housing costs are higher in the places where the supply evidence predicts they should be.
I once again focus on areas near ward boundaries using a border-discontinuity design.
Listed rents are about 3\% higher on the more stringent side of a ward border after controlling for listing characteristics and nearby amenities.
Home sale prices also appear to be slightly higher on the more stringent side, at about 1.5\%.
These results are consistent with lower supply on the more stringent side of ward boundaries and suggests that the uncertainty induced by informal regulations and political actors can have direct effects on affordability.

The paper contributes most directly to work on local control over housing supply and how incentives of particular neighborhoods and the aggregate city may not be aligned. 
\citet{khan_decentralized_2021} studies the same Chicago setting and finds that aldermen approve smaller rezonings near ward boundaries where benefits spill over to adjacent wards.
I study a broader margin by examining realized new construction rather than only statutory changes in allowed density through rezoned parcels to measure how the built environment can change in response to the regulatory environment.

The paper also relates to work on permitting, time-to-build, and fragmented land-use authority.
\citet{mayer_land_2000} and \citet{mayer_residential_2000} show that development lags are a first-order determinant of the supply response to demand shocks.
\citet{paciorek_supply_2013} shows that regulation reduces housing supply primarily through its effect on new construction timelines.
\citet{soltas_how_2026} provides within-parcel evidence that permitting costs are large and that developers value lots that have permits ``ready-to-issue (RTI)''.
This paper shows that the informal layer of the regulatory process and the individual decisionmakers involved can impact housing supply as well.

I also connect to the literature on how the fragmentation of decision-making authority leads to undersupply of housing.
\citet{fischel_homevoter_2009} argues that homeowners use local political institutions to restrict development and protect property values.
Subsequent work has documented how this plays out across municipalities, with each jurisdiction restricting development to a degree that may be locally rational but is collectively costly \citep{hilber_origins_2013, ortalo-magne_political_2014, kulka_how_2022, monarrez_dividing_2023, mast_warding_2024, bordeu_commuting_2025}.
My results show that a similar coordination failure operates within a single city when informal authority is fragmented across local representatives, who can create substantial variation in the regulatory environment that changes discretely at political boundaries.

The rest of the paper proceeds as follows.
Section~\ref{sec:data} describes the institutional setting and data.
Section~\ref{sec:model} presents a stylized time-to-build framework that organizes the supply predictions.
Section~\ref{sec:empirical_results} presents the main results.
Section~\ref{sec:discussion} discusses interpretation and policy implications, and Section~\ref{sec:conclusion} concludes.
